% ============================================================
% ANNEXES
% ============================================================
\chapter{Annexes}
\label{ch:annexes}

\section{Prompts Few-Shot complets}

\subsection{Prompt système}

Le prompt système fixe le rôle du LLM (correcteur expert BAC béninois), impose une chaîne de pensée obligatoire en 6 étapes, définit le barème /3 avec pénalité de subjectivité (-1pt), et spécifie le format de sortie JSON.

\lstinputlisting[language=Python,caption=Prompt système AZIZA (extrait de fond.py)]{../aziza/pipeline/fond.py}

\subsection{Exemples Few-Shot}

Cinq exemples sont fournis au LLM, couvrant les cas :
\begin{itemize}
    \item Contresens majeur (score 0) : corruption limitée au secteur public.
    \item Contraction parfaite (score 3) : langues maternelles.
    \item Omissions (score 2) : trafic d'enfants, causes secondaires.
    \item Contresens avec thèse partielle (score 2) : militarisation des scrutins.
    \item Subjectivité détectée (score 2 -1) : opinion personnelle dans la contraction.
\end{itemize}

\section{Cahier des charges complet}

Le cahier des charges initial du projet est disponible dans le dépôt GitHub à la racine du projet (fichier \texttt{Aziza\_cdc\_ia\_pure\_contraction.md}). Il détaille :

\begin{itemize}
    \item Le contexte et la justification du projet.
    \item La problématique et les objectifs (général et spécifiques).
    \item L'architecture technique du pipeline IA.
    \item Le plan d'expérimentation et les métriques d'évaluation.
    \item Les livrables et le calendrier.
\end{itemize}

\section{Extraits de code}

\subsection{Orchestrateur du pipeline}

\lstinputlisting[language=Python,caption=Pipeline complet (complet.py)]{../aziza/pipeline/complet.py}

\subsection{Machine à états de l'apprenant}

\lstinputlisting[language=Python,caption=Machine à états (machine_etats.py)]{../aziza/machine_etats.py}

\section{Captures d'écran de l'interface}

\subsection{Swagger UI}

L'interface Swagger est accessible à l'adresse \url{https://Jiovany11.pythonanywhere.com/apidocs/} et documente les 8 endpoints de l'API REST \AZIZA.

\subsection{Page de résultats}

La page de résultats affiche le breakdown complet : scores Forme (/7), Fond (/3), verrous, feedback pédagogique, et progression de l'apprenant.

\section{Structure des dépendances}

\begin{lstlisting}[caption=Dépendances Python (requirements.txt)]
flask>=3.0,<4.0
scikit-learn>=1.3
numpy>=2.0
requests>=2.31
language-tool-python>=2.7.1
python-dotenv>=1.0
flasgger>=0.9.7
flask-limiter>=3.8
\end{lstlisting}
